% =====================================================================
% Stage 2 -- High-Level Refinement.
% =====================================================================

\begin{requirement}{FRE-COMPAT.UNSIZED-01}[draft][functionalRequirement]{Bodies of unknown length}

  \localreqdiagram*{rfrunsized01}{\reqid}{\reqdesc}

  \texttt{pingora-cache} hands the backend a response body as it arrives from
  the upstream: the miss handler obtained from \texttt{get\_miss\_handler}
  receives the body chunk by chunk, and the end of the body is signalled only
  with the final chunk. Nothing in the \texttt{Storage} contract declares the
  body's total size when the store opens --- a chunked upstream response
  carries no \texttt{Content-Length} to declare.

  WHEN a store begins without the body's total length being declared, the
  backend shall persist the complete body while holding in process memory no
  more of that body than the chunk its current \texttt{write\_body} call is
  writing, independent of the body's length.

  The bound rules out buffering a whole body before writing it, and any
  per-store buffer besides: the operating system's page cache already
  batches writes to the device, so a buffer in the backend would save only
  system calls, at the cost of memory per concurrent store. The chunk being
  written is allowed because \texttt{pingora-cache} has already allocated
  it when the backend receives it, and its size depends on the
  upstream protocol --- HTTP/1 bodies
  are read through a 64~KiB buffer,\footnote{\pgsrc{pingora-core/src/protocols/http/v1/body.rs}{31}{\texttt{BODY\_BUFFER\_SIZE}}
  is \texttt{1024 * 64}.} while an HTTP/2 data frame can be larger.

  \textit{Basis: \texttt{Storage::get\_miss\_handler} and the
  \texttt{MissHandler} chunked write path.}\footnote{\pgsrc{pingora-cache/src/storage.rs}{176}{\texttt{Storage::get\_miss\_handler}}
  returns a \pgsrc{pingora-cache/src/storage.rs}{399}{\texttt{MissHandler}},
  whose \pgsrc{pingora-cache/src/storage.rs}{357}{\texttt{HandleMiss::write\_body(data,~eof)}}
  is called once per chunk, with \texttt{eof} set only on the final call.}

  \reqdiaglink{overview}{\textit{View the domain overview.}}

  \reqcover{verify}{TC}

  \deriveReqt{HLR-COMPAT.PINGORA-01}
\end{requirement}

\begin{requirement}{FRE-COMPAT.CONCURRENT-01}[draft][functionalRequirement]{Concurrent same-key writes}

  \localreqdiagram*{rfrconcurrent01}{\reqid}{\reqdesc}

  \texttt{pingora-cache} permits several writes of the same key to be in
  flight at once. The \texttt{CacheLock} that serialises miss handling is an
  optional collaborator, and even when configured its age timeout deliberately
  lets a new writer start while a previous one may still be writing; during a
  hot upgrade, a peer process holds its own, independent lock table
  (\reqref{FRE-COMPAT.HOTUPGRADE-01}), so same-key writes can also race
  across processes.

  While several writes of the same key are in flight at once, the backend
  shall keep each in isolation --- no writer's bytes ever interleaved with
  another's --- and shall make each write \emph{durably published} ---
  visible by name to a scan, a restart, or a peer process --- as a complete
  object or not at all. This is the general form of the atomicity guarantee
  whose replacement-specific case is detailed in
  \reqref{FRE-COMPLIANCE.REPLACE-01}; a lookup within the writing process
  may additionally be served a partial hit of an in-flight write, as
  governed by \reqref{FRE-COMPAT.STREAMING-01}.
  When those writes complete, the backend shall converge on a single winner;
  the atomic-replacement semantics of that convergence are also owned by
  \reqref{FRE-COMPLIANCE.REPLACE-01}.

  \textit{Basis: \texttt{HttpCache}'s optional cache lock and its age
  timeout;}\footnote{Once a lock ages out or is abandoned,
  \pgsrc{pingora-cache/src/lock.rs}{110}{\texttt{CacheLock::lock}} stops
  treating it as an in-progress write and hands out a fresh write permit
  instead of making the caller wait --- the mechanism behind the age timeout
  described above.} \textit{the stock \texttt{MemCache}'s per-writer keying
  of in-flight
  writes.}\footnote{\pgsrc{pingora-cache/src/memory.rs}{364}{\texttt{MemCache::get\_miss\_handler}}
  assigns each new write its own incrementing id and inserts it into the
  per-key map alongside any other writes already in flight.}

  \reqdiaglink{overview}{\textit{View the domain overview.}}

  \reqcover{verify}{TC}

  \deriveReqt{HLR-COMPAT.PINGORA-01}
  \trace{FRE-COMPLIANCE.REPLACE-01}
  \trace{FRE-COMPAT.HOTUPGRADE-01}
  \trace{FRE-COMPAT.STREAMING-01}
\end{requirement}

\begin{requirement}{FRE-COMPAT.EVICTION-01}[draft][functionalRequirement]{Eviction state consistency}

  \localreqdiagram*{rfrcompat01}{\reqid}{\reqdesc}

  \texttt{pingora-cache} manages capacity through an eviction manager (the
  \texttt{EvictionManager} trait, e.g.\ its \texttt{simple\_lru} and \texttt{lru}
  implementations) that holds, in memory, the set of cached items, their on-disk
  sizes, and the recency/freshness ordering used to choose eviction victims. The
  eviction manager is a component \emph{distinct} from the \texttt{Storage}
  backend, and it is the framework --- not the backend --- that drives it: both
  objects are handed to \texttt{HttpCache} (e.g.\ through the cache's
  \texttt{enable} configuration set up in the proxy), after which
  \texttt{HttpCache} records admissions and accesses against the manager as objects
  are written and read, and purges the victims it returns through the backend's
  removal path. The two are configured independently and the backend never calls
  the manager itself.

  Because this backend persists cached objects across process restarts
  (\reqref{HLR-PERSISTENCE.DURABLE-01}) while the
  eviction manager's in-memory state does not survive on its own, a naive restart
  leaves capacity accounting and eviction decisions out of step with the objects
  actually persisted: objects the manager no longer knows about are never counted
  against the size budget and never evicted. The discrepancy exists \emph{only at
  restart} --- in steady state \texttt{HttpCache} keeps a single manager in step ---
  so the requirement is to close that startup gap.

  A direct consequence of the framework owning the runtime \texttt{admit},
  \texttt{access}, and purge calls is that \emph{a \texttt{Storage}
  implementation cannot, on its own, keep the eviction manager synchronised}:
  its own contribution is to persist objects durably and to expose enumeration
  of the persisted entries (key, size, freshness), but bringing the manager
  back in step with that persisted set after a restart requires cooperation
  beyond what the \texttt{Storage} trait itself can provide. Meeting this requirement therefore demands more than a
  pure drop-in \texttt{Storage} backend --- a compatibility consideration for
  \reqref{HLR-COMPAT.PINGORA-01}. This requirement thus derives from both
  parents: \reqref{HLR-PERSISTENCE.DURABLE-01} creates the gap by making
  objects outlive the manager's memory, and \reqref{HLR-COMPAT.PINGORA-01}
  fixes the framework-owned eviction contract the gap must be closed within.

  WHEN the backend reports its restart procedure complete, the eviction
  manager's view of admitted objects shall be consistent with the durably
  published objects: every durably published object represented once ---
  keyed by its \texttt{CompactCacheKey}, with its record size and stored
  freshness deadline --- no entry remaining for an object that is no longer
  present, and the manager's reported \texttt{total\_items} and
  \texttt{total\_size} equal to the count and aggregate record size of those
  objects.

  An object's \emph{record size} is the full length of the file that stores
  it, metadata and body together. The backend reports that same size both
  when a store completes and when the restart procedure admits the object, so
  the two admission paths agree and \texttt{total\_size} measures the disk
  the cache occupies. Two things are deliberately left open. Before the
  restart procedure completes --- serving may already have begun ---
  capacity accounting may be inaccurate. And the recency order the manager
  used to choose victims before the restart need not survive it: afterwards,
  victims may be chosen in a different order.

  \textit{Basis: \texttt{HttpCache} drives the eviction manager on the hit
  and miss paths;}\footnote{\pgsrc{pingora-cache/src/lib.rs}{880}{\texttt{eviction.access(...)}}
  is called on a hit and
  \pgsrc{pingora-cache/src/lib.rs}{1118}{\texttt{eviction.admit(...)}} on a
  completed miss.} \textit{the \texttt{EvictionManager} trait's
  \texttt{admit}/\texttt{access}
  contract;}\footnote{\pgsrc{pingora-cache/src/eviction/mod.rs}{230}{\texttt{trait EvictionManager}}
  declares \pgsrc{pingora-cache/src/eviction/mod.rs}{252}{\texttt{admit}} and
  \pgsrc{pingora-cache/src/eviction/mod.rs}{289}{\texttt{access}}.} \textit{item
  identity via
  \texttt{CacheEntryKey};}\footnote{\pgsrc{pingora-cache/src/eviction/mod.rs}{53}{\texttt{enum CacheEntryKey}}
  wraps a \pgsrc{pingora-cache/src/key.rs}{153}{\texttt{struct CompactCacheKey}},
  optionally paired with a storage-defined \texttt{CacheEntryId}. A backend that
  publishes a single object per key supplies no identifier
  (\reqref{ADR-COMPAT.ENTRYID-01}), so the entry key is \texttt{KeyOnly} and
  carries the cache key alone.}
  \textit{the manager's aggregate
  accounting.}\footnote{\pgsrc{pingora-cache/src/eviction/mod.rs}{232}{\texttt{total\_size}}
  and \pgsrc{pingora-cache/src/eviction/mod.rs}{234}{\texttt{total\_items}}
  report the manager's tracked byte size and item count.}

  \reqdiaglink{overview}{\textit{View the domain overview.}}

  \reqcover{verify}{TC}
  \reqcover{satisfy}{ADR} % An architecture decision selecting the reconstruction strategy

  \deriveReqt{HLR-COMPAT.PINGORA-01}
  \deriveReqt{HLR-PERSISTENCE.DURABLE-01}
\end{requirement}

\begin{requirement}{FRE-COMPAT.HOTUPGRADE-01}[draft][functionalRequirement]{Hot-upgrade cache consistency}

  \localreqdiagram*{rfrhotupgrade01}{\reqid}{\reqdesc}

  Pingora supports \emph{hot} (graceful, zero-downtime) upgrades: a new
  (\emph{incoming}) executable takes over the listening sockets from the
  running (\emph{outgoing}) one, which drains its already-accepted
  in-flight requests before exiting, so for a bounded window two processes
  run concurrently. Because this backend persists objects durably
  (\reqref{HLR-PERSISTENCE.DURABLE-01}), the persisted state outlives any
  single process --- so both executables, configured for the same
  upstreams, access the very same cache during that window, a concern that
  would not arise if state were confined to one process's memory. Only
  objects already durably published are shared between the two processes
  --- in-flight activity such as streaming partial hits
  (\reqref{FRE-COMPAT.STREAMING-01}) is not yet part of the cache, and is
  outside this requirement's scope.

  While two processes share the same durable cache during a hot upgrade,
  the backend in each shall keep the shared cache consistent:
  \begin{itemize}
    \item \textbf{Outgoing process.} WHEN the outgoing process receives
      the Pingora upgrade signal, its backend shall make no further durable
      change to the shared cache: a store that completes afterwards shall
      not be durably published, and a purge or metadata update requested
      afterwards shall fail with an error rather than be reported as
      applied. It shall keep serving reads of objects already durably
      published, by either process, so that its remaining requests
      complete.
    \item \textbf{Incoming process.} Its backend is subject to the ordinary
      concurrent-writer and eviction-consistency guarantees this backend
      already provides for a single running instance
      (\reqref{FRE-COMPAT.CONCURRENT-01}, \reqref{FRE-COMPAT.EVICTION-01})
      as though it were the sole writer --- hot upgrade grants no exception
      to those --- and none of its storage operations, the restart-time
      reclaim scan included, shall make an object unreadable through a file
      descriptor the outgoing process already holds open.
  \end{itemize}

  A purge refused by the draining process surfaces as a failed purge
  request, which its client can retry against the incoming process; reporting
  it as applied would break the purge-acknowledgement guarantee
  (\reqref{FRE-COMPLIANCE.PURGEDURABLE-01}). How the outgoing process's backend
  learns that the signal has arrived is left to the design: the server
  announces the transition to its own code through an execution-phase
  broadcast.\footnote{\pgsrc{pingora-core/src/server/mod.rs}{235}{\texttt{Server::watch\_execution\_phase}}
  subscribes to the phase broadcast, which announces
  \texttt{ExecutionPhase::GracefulUpgradeTransferringFds} when the upgrade
  signal (\texttt{SIGQUIT}) arrives
  (\pgsrc{pingora-core/src/server/mod.rs}{279}{L279}); both processes then
  accept connections for \texttt{CLOSE\_TIMEOUT}
  (\pgsrc{pingora-core/src/server/mod.rs}{66}{5~s}) before the old one stops
  listening and drains.}

  This requirement derives from both parents: \reqref{HLR-COMPAT.PINGORA-01}
  brings the Pingora hot-upgrade model, with its window of two processes, and
  \reqref{HLR-PERSISTENCE.DURABLE-01} makes the cache state they share
  outlive either process.

  \textit{Basis: Pingora graceful (zero-downtime) upgrade
  (\reqref{HLR-COMPAT.PINGORA-01}).}

  \reqdiaglink{overview}{\textit{View the domain overview.}}

  \reqcover{verify}{TC}
  \reqexpect{satisfy}{FDS-EVICTION.STARTUP-01}{Realise the reconstruction
  procedure so the incoming process's eviction view is rebuilt from durably
  published objects only, keeping the shared cache consistent across the
  hot-upgrade window.}
  \reqexpect{satisfy}{FDS-COMPAT.READONLY-01}{Specify how the outgoing process's
  backend learns that the upgrade signal arrived, what it suppresses from
  then on (publication, purge, metadata update) and how it reports refusals,
  and how the incoming process's reconstruction is ordered relative to the
  outgoing process's last publications.}

  \deriveReqt{HLR-COMPAT.PINGORA-01}
  \deriveReqt{HLR-PERSISTENCE.DURABLE-01}
  \trace{FRE-COMPAT.CONCURRENT-01}
  \trace{FRE-COMPAT.EVICTION-01}
  \trace{FRE-COMPAT.STREAMING-01}
  \trace{FRE-COMPLIANCE.PURGEDURABLE-01}
\end{requirement}

\begin{requirement}{FRE-COMPAT.STREAMFLAG-01}[draft][functionalRequirement]{Advertise streaming writes}

  \localreqdiagram*{rfrstreamflag01}{\reqid}{\reqdesc}

  The backend shall report \texttt{true} from its implementation of the
  \texttt{Storage} trait's \texttt{support\_streaming\_partial\_write}
  method.

  The stampede protection in \texttt{pingora-cache} assumes a backend that can serve
  an object \emph{while it is still being written}: when \texttt{Storage}
  advertises \texttt{support\_streaming\_partial\_write}, \texttt{HttpCache}
  releases the cache lock as soon as the leader's miss handler opens --- right
  after \texttt{get\_miss\_handler} rather than at
  \texttt{finish\_miss\_handler} --- and every concurrent request for the key
  is then expected to be served by a \emph{partial} read of the in-flight
  object. A backend without partial reads instead holds every follower, under
  the lock, for the leader's whole body; for a large or slow upstream response
  that turns stampede protection into a head-of-line delay of the full
  download for every concurrent client. Advertising the flag is what lets
  followers reach the partial hits of \reqref{FRE-COMPAT.STREAMING-01} at all,
  so the flag is only truthful together with that requirement,
  \reqref{FRE-COMPAT.STREAMDELIVERY-01}, and \reqref{FRE-COMPAT.STREAMEND-01}.

  \textit{Basis: the \texttt{Storage} trait's
  \texttt{support\_streaming\_partial\_write}
  flag;}\footnote{\pgsrc{pingora-cache/src/storage.rs}{245}{\texttt{Storage::support\_streaming\_partial\_write}}
  defaults to \texttt{false}; a backend opts in by overriding it.}
  \textit{\texttt{HttpCache} releasing the cache lock right after the miss
  handler opens, instead of waiting for the store to
  finish.}\footnote{\pgsrc{pingora-cache/src/lib.rs}{997}{\texttt{HttpCache::set\_miss\_handler}}
  releases the lock immediately once
  \texttt{support\_streaming\_partial\_write()} is true, rather than at
  \texttt{finish\_miss\_handler}.}

  \reqdiaglink{overview}{\textit{View the domain overview.}}

  \reqcover{verify}{TC}
  \reqcover{satisfy}{FDS}

  \deriveReqt{HLR-COMPAT.PINGORA-01}
  \trace{FRE-COMPAT.STREAMING-01}
\end{requirement}

\begin{requirement}{FRE-COMPAT.STREAMING-01}[draft][functionalRequirement]{Partial hits on lookup}

  \localreqdiagram*{rfrstreaming01}{\reqid}{\reqdesc}

  WHILE one or more stores of a key that no invalidating purge has ended
  (\reqref{FRE-COMPLIANCE.PURGE-01}) are in flight --- each from the opening
  of its miss handler until it completes or aborts --- the backend shall
  answer a \texttt{lookup} of that key with a \emph{partial hit} of the most
  recently opened of those stores: that store's object metadata together with a
  \texttt{HitHandler} that streams the body bytes the store writes.

  The partial hit takes precedence over any published object of the same key,
  exactly as the stock \texttt{MemCache} prefers a partial read on lookup.
  \texttt{pingora-cache} permits several stores of one key to be in flight at
  once; resolving to the most recently opened one gives a follower the
  freshest response the cache is fetching. What the partial hit's reader
  delivers, and when, is stated in \reqref{FRE-COMPAT.STREAMDELIVERY-01} and
  \reqref{FRE-COMPAT.STREAMEND-01}; what it owes when the store aborts, in
  \reqref{FRE-COMPLIANCE.PARTIALABORT-01}.

  \textit{Basis: the stock \texttt{MemCache}'s unconditional preference for a
  partial read on
  lookup.}\footnote{\pgsrc{pingora-cache/src/memory.rs}{313}{\texttt{MemCache::lookup}}
  always prefers an in-flight write over a complete cached entry; per its own
  comment it has no preference among multiple in-flight writers, so the
  most-recently-opened tie-break above is this backend's own addition, not
  inherited framework behavior.}

  \reqdiaglink{overview}{\textit{View the domain overview.}}

  \reqcover{verify}{TC}
  \reqexpect{satisfy}{FDS-STORAGE.PARTIAL-01}{Specify the in-flight write
  registry, the lookup order (in-flight write before the published record),
  and the most-recently-opened tie-break.}

  \deriveReqt{HLR-COMPAT.PINGORA-01}
\end{requirement}

\begin{requirement}{FRE-COMPAT.STREAMDELIVERY-01}[draft][performanceRequirement]{Partial-hit chunk delivery}

  \localreqdiagram*{rfrstreamdelivery01}{\reqid}{\reqdesc}

  WHILE a partial hit's reader is attached to an in-flight store, the reader
  shall deliver the body bytes that store had already written, followed by
  the chunks the store subsequently appends, with at least 99\,\% of those
  appended chunks delivered within 10~ms of their \texttt{write\_body} call
  returning.

  A chunk is delivered when the reader's \texttt{read\_body} call returns it;
  the bound applies on a host that is not CPU-saturated. The follower thus
  streams the body as the leader writes it, without waiting for the store to
  finish. The bound is a percentile rather than a worst case so that a single
  scheduling delay does not violate it, while still excluding designs that
  poll the file or batch notifications, whose delays are typically an order
  of magnitude longer.

  \reqdiaglink{overview}{\textit{View the domain overview.}}

  \reqcover{verify}{TC}
  \reqexpect{satisfy}{FDS-STORAGE.PARTIAL-01}{Specify how a reader learns of
  each appended chunk so that it can deliver the chunk within the latency
  bound.}

  \deriveReqt{HLR-COMPAT.PINGORA-01}
  \trace{FRE-COMPAT.STREAMING-01}
\end{requirement}

\begin{requirement}{FRE-COMPAT.STREAMEND-01}[draft][functionalRequirement]{Partial-hit completion}

  \localreqdiagram*{rfrstreamend01}{\reqid}{\reqdesc}

  WHEN the in-flight store a partial hit's reader is attached to completes,
  the reader shall deliver the store's remaining body bytes and then end the
  body normally, whether or not the object has meanwhile been durably
  published.

  Completion is followed by the object's durable publication, which may
  happen while a follower is still reading, so a follower's normal end cannot
  rely on the in-flight store still existing when it gets there. When the store instead
  aborts --- an upstream failure or a cancellation --- the effect on an
  attached reader is governed by
  \reqref{FRE-COMPLIANCE.PARTIALABORT-01}.

  \reqdiaglink{overview}{\textit{View the domain overview.}}

  \reqcover{verify}{TC}
  \reqexpect{satisfy}{FDS-STORAGE.PARTIAL-01}{Specify how each reader
  delivers the remainder and ends normally on completion, and why its
  reader stays valid across the object's durable publication.}

  \deriveReqt{HLR-COMPAT.PINGORA-01}
  \trace{FRE-COMPAT.STREAMING-01}
  \trace{FRE-COMPLIANCE.PARTIALABORT-01}
\end{requirement}

\begin{requirement}{FRE-COMPAT.FORMATUPGRADE-01}[draft][functionalRequirement]{Records across Pingora changes}

  \localreqdiagram*{rfrformatupgrade01}{\reqid}{\reqdesc}

  The metadata the backend persists is the serialisation, by
  \texttt{pingora-cache} itself, of each object's
  \texttt{CacheMeta},\footnote{\pgsrc{pingora-cache/src/meta.rs}{890}{\texttt{CacheMeta::serialize}}
  yields the internal-metadata and header byte segments the backend stores.}
  and that encoding is not the backend's to control. Its internal-metadata
  segment carries a version the framework decodes across its own
  releases,\footnote{\pgsrc{pingora-cache/src/meta.rs}{166}{\texttt{internal\_meta::deserialize}}
  decodes versions 0--2 and rejects any other
  (\pgsrc{pingora-cache/src/meta.rs}{190}{L190}).} and its header segment is
  compressed with a dictionary each process loads at
  startup.\footnote{\pgsrc{pingora-cache/src/meta.rs}{965}{\texttt{HEADER\_SERDE}}
  is built from the dictionary set through
  \pgsrc{pingora-cache/src/meta.rs}{989}{\texttt{set\_compression\_dict\_path}}
  or \pgsrc{pingora-cache/src/meta.rs}{1008}{\texttt{set\_compression\_dict\_content}},
  if any.} Because persisted records outlive the process that wrote them
  (\reqref{HLR-PERSISTENCE.DURABLE-01}), a cache directory written by a
  process built against one \texttt{pingora-cache} release may be opened by a
  process built against another --- most visibly across a hot upgrade that
  also changes the Pingora version (\reqref{FRE-COMPAT.HOTUPGRADE-01}). This requirement derives from both
  parents: \reqref{HLR-COMPAT.PINGORA-01} fixes the framework whose encoding
  may change, and \reqref{HLR-PERSISTENCE.DURABLE-01} makes records survive
  such a change.

  WHEN the backend reads a record written by a process that used a different
  \texttt{pingora-cache} release, the backend shall either return the object
  exactly as that record stored it or report a miss for it, and shall never
  return incorrectly decoded metadata.

  The header compression dictionary is outside this requirement.
  \texttt{pingora-cache} offers no way to learn which dictionary a process
  loaded, nor which one a stored header segment was compressed with, so the
  backend cannot tell a record written under another dictionary from one
  written under its own. Every process sharing a cache directory is
  therefore expected to load the same dictionary, and a deployment that
  changes it deletes the cache directory or gives the new process a new one.
  This is a deployment constraint the deliverable documents, like the
  requirement that processes sharing a directory run on one host
  (\reqref{ADR-STORAGE.WRITERID-01}).

  The backend meets this without tying its records to any Pingora release.
  Each record carries the backend's \emph{own} format version
  (\reqref{SDS-LAYOUT.OBJECTFILE-01}), which the backend changes whenever
  records written before would no longer be read back exactly: a change to
  its own layout, or a \texttt{pingora-cache} release whose encoding of the
  metadata differs. A record of another version is reported as a miss. A
  \texttt{pingora-cache} release that leaves the encoding alone therefore
  keeps the cache warm; one that changes it starts the cache cold once.

  Whether a release changed the encoding is a decision the backend's
  maintainers make, not an inference from a version number. The Pingora
  release numbering cannot settle it: 0.9.0 extended \texttt{InternalMetaV2} with new
  fields while keeping its version byte at 2, compatibly, so 0.8.1's metadata
  would have decoded under 0.9.0, yet a release could as well change the
  encoding under an unchanged version byte.
  \reqref{TC-COMPAT.FORMATUPGRADE-01} holds the encoding the current format
  relies on as fixed bytes, and fails when the release in use reads or
  writes them differently, so no release is adopted without that decision.
  The guarantee is therefore verified for every release that test has run
  against. A release an integrator resolves beyond those is assumed, not
  verified, to keep the encoding, as a patch release of the 0.9 line
  should.

  \reqdiaglink{overview}{\textit{View the domain overview.}}

  \reqcover{satisfy}{SDS}
  \reqcover{verify}{TC}

  \deriveReqt{HLR-COMPAT.PINGORA-01}
  \deriveReqt{HLR-PERSISTENCE.DURABLE-01}
  \trace{FRE-COMPAT.HOTUPGRADE-01}
\end{requirement}

\begin{requirement}{FRE-COMPAT.WRITEFAIL-01}[draft][functionalRequirement]{Write failures fail the store}

  \localreqdiagram*{rfrwritefail01}{\reqid}{\reqdesc}

  IF writing a store's data to persistent storage fails --- for example
  because the file system is full --- THEN the backend shall return an error
  from the \texttt{write\_body} or \texttt{finish} call during which the
  failure occurred.

  The \texttt{HandleMiss} contract gives each call a \texttt{Result} for
  exactly this purpose.\footnote{\pgsrc{pingora-cache/src/storage.rs}{357}{\texttt{HandleMiss::write\_body}}
  returns \texttt{Result<()>}.} A store that has failed this way is
  incomplete, so what lookups may then observe of it is governed by
  \reqref{FRE-COMPLIANCE.INCOMPLETE-01}, and the storage it consumed by
  \reqref{FRE-COMPAT.ABORTSPACE-01}.

  \reqdiaglink{overview}{\textit{View the domain overview.}}

  \reqexpect{satisfy}{FDS-STORAGE.PUBLISH-01}{Specify how a failed write or
  sync during a store is detected and surfaced as an error from the call in
  progress, and that the store is then treated as aborted.}
  \reqcover{verify}{TC}

  \deriveReqt{HLR-COMPAT.PINGORA-01}
  \trace{FRE-COMPLIANCE.INCOMPLETE-01}
  \trace{FRE-COMPAT.ABORTSPACE-01}
\end{requirement}

\begin{requirement}{FRE-COMPAT.ABORTSPACE-01}[draft][functionalRequirement]{Aborted stores release space}

  \localreqdiagram*{rfrabortspace01}{\reqid}{\reqdesc}

  WHEN a store ends without being durably published --- aborted by an
  upstream failure, a cancellation, or a write failure --- the backend shall
  remove the store's temporary no later than when its miss handler is
  dropped, and shall free the storage that store consumed without waiting for
  a restart.

  Space held by an unpublished store is not counted by the eviction manager,
  which admits an object only once its store completes, so a store that
  failed without releasing its space would shrink the disk available to the
  cache with every upstream failure until the next restart. Freeing a
  large file's storage can take a long time, so it may complete
  shortly after the drop, away from the thread that dropped the handler; a
  reader still streaming the store keeps the storage until it finishes
  (\reqref{FDS-STORAGE.PARTIAL-01}). The stock
  backend releases an unfinished write the same
  way.\footnote{\pgsrc{pingora-cache/src/memory.rs}{291}{\texttt{Drop for MemMissHandler}}
  removes the handler's in-flight write.} A process crash is different ---
  no handler is dropped --- and its leftovers are reclaimed as
  \reqref{FRE-COMPAT.STARTRECLAIM-01} requires.

  \reqdiaglink{overview}{\textit{View the domain overview.}}

  \reqexpect{satisfy}{FDS-STORAGE.PUBLISH-01}{Specify that an unpublished
  store's per-writer temporary is unlinked when the store ends, including on
  abort and on write failure, and how attached partial-hit readers are
  unaffected by that unlink.}
  \reqcover{verify}{TC}

  \deriveReqt{HLR-EFFICIENCY.SPACE-01}
  \trace{FRE-COMPLIANCE.INCOMPLETE-01}
  \trace{FRE-COMPAT.WRITEFAIL-01}
  \trace{FRE-COMPAT.STARTRECLAIM-01}
\end{requirement}

\begin{requirement}{FRE-COMPAT.STARTRECLAIM-01}[draft][functionalRequirement]{Reclaim leftovers at start}

  \localreqdiagram*{rfrstartreclaim01}{\reqid}{\reqdesc}

  WHEN a process of the backend starts, the backend shall release the storage
  held by unpublished stores whose writing process is no longer running.

  Such a leftover is not reachable through a keyed \texttt{Storage} call
  again: a process crash drops no miss handler, so the release on abort
  (\reqref{FRE-COMPAT.ABORTSPACE-01}) never runs for its stores. Without a
  release at start, each crash would shrink the disk available to the cache
  for good. A store whose writer is still running --- such as a hot-upgrade
  peer's --- is not a leftover, and the incoming process leaves it untouched
  (\reqref{FRE-COMPAT.HOTUPGRADE-01}).

  A \texttt{Vary} variant record whose primary key holds no stored object is
  not a leftover either. It is a published object, which the eviction manager
  counts against the cache's size budget and evicts like any other
  (\reqref{FRE-COMPAT.EVICTION-01}), and a lookup may still reach it: an
  integrator's cache key can carry a variance of its own, and no primary is
  then ever stored.

  \reqdiaglink{overview}{\textit{View the domain overview.}}

  \reqcover{satisfy}{FDS}
  \reqcover{verify}{TC}

  \deriveReqt{HLR-EFFICIENCY.SPACE-01}
  \trace{FRE-COMPLIANCE.INCOMPLETE-01}
  \trace{FRE-COMPAT.ABORTSPACE-01}
\end{requirement}

\begin{requirement}{FRE-COMPAT.EXPIRE-01}[draft][functionalRequirement]{Expire keeps the body}

  \localreqdiagram*{rfrexpire01}{\reqid}{\reqdesc}

  WHEN the framework asks the storage backend to expire a stored object, the
  backend shall leave that object in the stale state.

  An object is in the \emph{stale state} when a later lookup for its key
  returns it, with its stored body readable in full, and reports it as not
  fresh at any instant after the expiry. Deleting the object is
  therefore not the stale state, and neither is leaving it fresh.

  The boundary instant itself is deliberately excluded, because the framework
  owns both halves of that comparison and places it there.
  \texttt{CacheMeta::expire\_at(instant)}, which Pingora documents as marking
  an asset ``stale as of
  \emph{instant}'',\footnote{\pgsrc{pingora-cache/src/meta.rs}{800}{\texttt{CacheMeta::expire\_at}}
  sets \texttt{fresh\_until} to the earlier of its current value and
  \emph{instant}.} lowers the stored freshness deadline to that instant, while
  \texttt{CacheMeta::is\_fresh} compares
  inclusively.\footnote{\pgsrc{pingora-cache/src/meta.rs}{716}{\texttt{CacheMeta::is\_fresh}}
  returns \texttt{fresh\_until >= time}.} An object expired at $t$ is therefore
  reported fresh at exactly $t$ and stale at every instant after it --- a
  distinction below the one-second resolution RFC~9111 works at, and one no
  backend storing the deadline the framework derives can alter.

  Pingora 0.9.0 added this second removal operation beside
  \texttt{purge}.\footnote{\pgsrc{pingora-cache/src/storage.rs}{226}{\texttt{Storage::expire}}.}
  Where a purge deletes the object and costs the next request a full refetch,
  an expiry keeps it and costs only a conditional request the origin can answer
  with a \texttt{304}, after which the retained body is served again. A proxy
  chooses between the two per request.\footnote{\pgsrc{pingora-proxy/src/proxy_trait.rs}{732}{\texttt{ProxyHttp::purge\_action}}.}

  The obligation is about the object's observable state, not about any
  particular field or file the backend rewrites to reach it. Two things it
  deliberately leaves alone. Whether a reader is served the retained body \emph{while} that
  revalidation is in flight is decided by the object's own serve-stale windows,
  which an expiry does not touch. And a store already in flight for the same key
  is not cancelled: if it completes it publishes a fresh object. An
  invalidating purge, by contrast, ends such a store
  (\reqref{FRE-COMPLIANCE.PURGE-01}), since it may carry the content the
  purge removes; an expiry asks only for a revalidation, which a fresh
  object provides.

  The trait supplies a default that deletes the object
  instead,\footnote{\pgsrc{pingora-cache/src/storage.rs}{230}{The default}
  forwards to \texttt{Storage::purge}.} which does keep the object from being
  served as fresh, but by removing it rather than by reaching the stale state,
  so it does not satisfy this requirement. What an expiry of a key holding
  nothing reports follows \reqref{FRE-COMPLIANCE.PURGEABSENT-01}.

  \reqdiaglink{overview}{\textit{View the domain overview.}}

  \reqcover{satisfy}{ADR}
  \reqcover{verify}{TC}

  \deriveReqt{HLR-COMPAT.PINGORA-01}
  \trace{FRE-COMPLIANCE.PURGE-01}
  \trace{FRE-COMPLIANCE.REVALIDATION-01}
\end{requirement}

\begin{requirement}{ADR-PERSISTENCE.MECHANISM-01}[draft][designConstraint]{Plain file system storage}

  \localreqdiagram*{adr1}{\reqid}{\reqdesc}

  \textbf{Context:} The cache must persist HTTP response headers and bodies
  durably (\reqref{HLR-PERSISTENCE.DURABLE-01}), this is a first version
  that favours the simplest adequate design (\reqref{HLR-SCOPE.V1-01}),
  and a body arrives as a stream whose total length is unknown when
  writing begins (\reqref{FRE-COMPAT.UNSIZED-01}). The choice is
  twofold: which storage substrate to build on, and whether to adopt an
  established storage engine or write the storage layer ourselves.

  \textbf{Decision:} Persist cached objects on a plain file system ---
  growable files in directories --- rather than adopting an embeddable
  storage engine or building on a raw block device, a custom preallocated
  arena, or an embedded database. The kernel's file system already does an
  allocator's work: block allocation and free-space management, growing
  files in place, keeping \emph{its own} structures intact across crashes
  (journaling), and giving every process on the host one coherent view of
  the same files. What it deliberately does not give, the backend keeps as
  its own duty:
  \begin{itemize}
    \item \emph{Atomic publication.} A file being written when the process
      dies is not rolled back, so object publication must be made atomic
      and leftovers discarded (\reqref{FRE-COMPLIANCE.INCOMPLETE-01},
      \reqref{FRE-COMPAT.STARTRECLAIM-01},
      \reqref{ADR-STORAGE.ATOMICPUBLISH-01}).
    \item \emph{Crash safety of its own writes.} A write can be reported as
      done while the bytes are still only in the operating system's cache,
      and syncing a file does not persist the directory entry that names
      it, so a host crash can lose or reorder writes. Which writes the
      backend synchronises, and what it detects instead, is its own
      decision (\reqref{ADR-PERSISTENCE.DURABILITY-01}). The eviction
      manager's state is not among them: it is rebuilt from the object
      records themselves at every start (\reqref{ADR-COMPAT.EVICTION-01}).
    \item \emph{Hot-upgrade coordination.} During a hot upgrade two
      processes see the same bytes but are not coordinated by the file
      system; making the outgoing process stop changing the shared cache,
      while leaving its open readers intact, is the backend's own protocol
      (\reqref{FRE-COMPAT.HOTUPGRADE-01}).
  \end{itemize}
  Every alternative below leaves those same coordination duties open ---
  or solves them at the price of a poorer fit for streaming bodies and a
  larger integration surface --- which is why the simplest substrate wins
  for a first version (\reqref{HLR-SCOPE.V1-01}). Appendix~\ref{app:alternatives}
  examines each alternative in detail.

  \textbf{Alternatives considered:}
  \begin{itemize}
    \item \emph{Adopt Apache Traffic Server's object store}: the most
      proven raw-disk HTTP cache, but not an adoptable component --- the
      store is built into Traffic Server's own process model and not
      published as a library,\footnote{The cache implementation lives
      under \texttt{src/iocore/cache}, part of IOCore (the process's core
      I/O and event system), and is built as an internal library linked
      only into the \texttt{traffic\_server} binary --- there is no
      separately packaged cache library.
      \url{https://github.com/apache/trafficserver/tree/10.1.4/src/iocore/cache}}
      so adopting it means carrying the server along.
    \item \emph{Adopt Meta's CacheLib}, an embeddable C++ caching engine:
      proven allocation, admission, and eviction
      machinery,\footnote{\url{https://cachelib.org/docs/Cache_Library_User_Guides/eviction_policy/\#eviction-policy}}
      and a body of unknown length fits as incrementally chained
      per-chunk
      allocations.\footnote{\url{https://cachelib.org/docs/Cache_Library_User_Guides/chained_items/\#chained-allocations}}
      Its eviction machinery need not conflict with
      the one in \texttt{pingora-cache}: the \texttt{EvictionManager}
      \texttt{pingora-cache} drives is
      optional,\footnote{\pgsrc{pingora-cache/src/lib.rs}{571}{\texttt{HttpCache::enable}}
      takes it as an \texttt{Option}; without it, nothing is evicted
      from storage.} so a backend can leave it unset and let CacheLib
      manage capacity on its own (\reqref{FRE-COMPAT.EVICTION-01}).
      Two costs rule it out for a first version instead:
      \begin{itemize}
        \item a C++ dependency, driven from Rust through an FFI boundary;
        \item one active cache instance per
          path,\footnote{\url{https://cachelib.org/docs/Cache_Library_User_Guides/Cache_persistence/\#common-failure-scenarios}}
          where a hot-upgrade window has two processes sharing one store
          (\reqref{FRE-COMPAT.HOTUPGRADE-01}).
      \end{itemize}
    \item \emph{Embedded database as the whole store} (libmdbx/LMDB,
      RocksDB, or pure-Rust engines such as \texttt{redb} and \texttt{fjall}):
      record packing and crash consistency come built in, and the pure-Rust
      engines remove the FFI cost as well, but their unit of storage is the
      record --- a body streaming in with unknown length
      (\reqref{FRE-COMPAT.UNSIZED-01}) must be chunked into records by the
      application --- and streaming partial hits
      (\reqref{FRE-COMPAT.STREAMING-01}) need the same in-memory
      coordination beside the store anyway. Cross-process use during a hot
      upgrade also varies by engine: RocksDB admits no multi-process
      writers,\footnote{\url{https://github.com/facebook/rocksdb/wiki/RocksDB-FAQ\#basic-readwrite}}
      so the incoming process could not write while the outgoing one still
      holds the database open, whereas the LMDB family allows several
      processes, one write transaction at a
      time.\footnote{\url{https://github.com/erthink/libmdbx\#gotchas}}
    \item \emph{Adopt an existing Rust disk cache} such as \texttt{cacache}, a
      key- and content-addressed cache with integrity-checked
      content:\footnote{\href{https://docs.rs/cacache/latest/cacache/}{\texttt{cacache}
      crate documentation}: ``a Rust library for managing local key and
      content address caches''; a streaming \texttt{Writer}'s data is
      committed only by \texttt{commit()}.} proven, no FFI, and streaming
      writes are supported. Rejected for this version: content becomes
      visible only when a write commits, so streaming partial hits still need
      the backend's own in-flight coordination; content addressing separates
      each object's metadata from its body, so a purge or metadata update
      leaves content that some unreferenced-content collection must remove,
      which its documentation does not provide; and a cryptographic integrity
      hash is computed over every body.
    \item \emph{Build on a raw block device or custom arena} (the
      substrate beneath Traffic Server's cache, CacheLib's flash engine,
      and Squid's Rock store): eliminates all per-object file-system overhead,
      and the systems built on it prove the model serves HTTP caching well.
      Choosing this substrate means building space
      management and an on-device directory from scratch. That
      combination is the largest bespoke surface of any alternative
      here, and the antithesis of a first version
      (\reqref{HLR-SCOPE.V1-01}).
  \end{itemize}

  \textbf{Consequences:}
  \begin{itemize}
    \item[+] Total cache size is bounded by disk capacity, not process RAM
      (\reqref{HLR-PERSISTENCE.CAPACITY-01}).
    \item[+] A growable file holds a body of unknown length with no size
      declared up front; bytes are appended as they stream in --- no
      per-chunk allocation and chaining, no pre-declared record sizes.
    \item[+] No foreign engine to integrate: no FFI boundary, no second
      eviction policy beside the one in \texttt{pingora-cache}, no engine-specific
      recovery tooling; the custom surface is confined to layout, naming,
      and the coordination duties above.
    \item[--] Ties the cache to file-system semantics: block-granular
      allocation with per-file overhead
      (\reqref{REQ-EFFICIENCY.FRAGMENTATION-01}) and per-directory entry
      scaling (\reqref{REQ-EFFICIENCY.DIRSCALE-01}) --- stated as derived
      requirements that the detailed decisions must satisfy.
    \item[--] An established engine would come with packing and
      compaction already built and tested; giving that up means the
      derived requirements above only set the efficiency targets --- the
      detailed design still has to build the machinery that meets them.
  \end{itemize}

  \reqdiaglink{overview}{\textit{View the domain overview.}}

  \reqexpect{deriveReqt}{REQ-EFFICIENCY.FRAGMENTATION-01}{Bound the aggregate
  space overhead a plain file system imposes on objects smaller than one
  block --- block rounding and per-file inode/directory-entry costs --- so
  that it grows sub-linearly with the number of such objects.}
  \reqexpect{deriveReqt}{REQ-EFFICIENCY.DIRSCALE-01}{Bound the number of
  entries that accumulate in any one directory, up to a stated cache
  capacity, so object lookup and maintenance stay fast as the number of
  cached objects grows.}

  \reqcover{refine}{ADR}

  \satisfy{HLR-PERSISTENCE.DURABLE-01}
  \satisfy{HLR-PERSISTENCE.CAPACITY-01}
  \satisfy{HLR-SCOPE.V1-01}
  \satisfy{FRE-COMPAT.UNSIZED-01}
  \trace{FRE-COMPAT.HOTUPGRADE-01}
  \trace{FRE-COMPAT.EVICTION-01}
  \trace{FRE-COMPLIANCE.INCOMPLETE-01}
\end{requirement}

\begin{requirement}{REQ-EFFICIENCY.FRAGMENTATION-01}[draft][requirement]{Sub-linear small-object overhead}

  \localreqdiagram*{reqfrag01}{\reqid}{\reqdesc}

  For cached objects smaller than one file-system block, the aggregate space
  the on-disk layout consumes beyond those objects' own metadata and body
  bytes shall grow sub-linearly with the number of such objects.

  A plain file system (\reqref{ADR-PERSISTENCE.MECHANISM-01}) allocates space
  in whole blocks (commonly 4~KiB), so a stored file's tail can waste up to
  nearly a full block. It also charges every file a fixed cost for its inode
  and directory entry, regardless of the file's size. Across many small
  cached objects --- a busy HTTP cache may hold many --- these costs grow in
  step with the object count and can waste a significant fraction of the
  disk. This requirement derives from both parents: the stakeholder
  space-efficiency goal (\reqref{HLR-EFFICIENCY.SPACE-01}) is what the waste
  works against, and the plain-file-system decision is what generates it.
  Verification is by analysis of the layout, confirmed by measuring disk use
  at growing object counts. The first functional version does not meet this
  requirement: \reqref{ADR-EFFICIENCY.COMPACTION-01} defers the packing that
  would, and declares the conflict.

  \reqdiaglink{overview}{\textit{View the domain overview.}}

  \reqcover{satisfy}{ADR}

  \deriveReqt{HLR-EFFICIENCY.SPACE-01}
  \deriveReqt{ADR-PERSISTENCE.MECHANISM-01}
\end{requirement}

\begin{requirement}{REQ-EFFICIENCY.DIRSCALE-01}[draft][requirement]{Directory-entry scaling}

  \localreqdiagram*{reqdirscale01}{\reqid}{\reqdesc}

  For a cache holding up to $10^8$ objects whose cache-key hashes are
  uniformly distributed, the on-disk layout shall hold no more than 2,000
  entries in any one directory.

  This requirement derives from two parents. On a plain file system
  (\reqref{ADR-PERSISTENCE.MECHANISM-01}) every cached object is reached
  through a directory entry, and many file systems slow down on name lookup
  and insertion as a single directory's entry count grows; listings and
  housekeeping scans grow more expensive with it. Left unbounded, that growth
  would make lookups slower as the cache fills, which
  \reqref{HLR-PERFORMANCE.LOOKUP-01} rules out. The capacity and the bound are
  a design target, deliberately beyond the $10^7$ objects that requirement
  measures, not a limit on what the cache may hold: beyond $10^8$ objects the
  bound no longer applies, but nothing else in this specification stops
  working. Verification is by analysis of the expected
  entry distribution for the layout, together with a test of the fan-out the
  analysis assumes.

  \reqdiaglink{overview}{\textit{View the domain overview.}}

  \reqexpect{satisfy}{SDS-LAYOUT.SHARDING-01}{Lay out objects so that, for up
  to $10^8$ uniformly hashed objects, no directory holds more than 2,000
  entries.}
  \reqcover{verify}{TC}

  \deriveReqt{ADR-PERSISTENCE.MECHANISM-01}
  \deriveReqt{HLR-PERFORMANCE.LOOKUP-01}
\end{requirement}

